\section{Prime-power refinement, full-exponent growth and formal scope}
\label{sec-prime-power-compatibility}

The general totient budget has a particularly sharp form at a common
prime-power divisor of the two actual pure exponents. Let $a,b>0$ be
coprime, $c=a+b$, $M=a^2+ab+b^2$ and $F=M^2+ab c^2$.
For the full prime-power part outside $1\bmod s$, use the notation
$c_{\mathrm{bad}}(s)$ of the preceding section.

\begin{corollary}[The full prime-power modulus]
\label{ppg-full-modulus}
Suppose $p\ge5$ is prime, $k\ge1$, $n=p^k$, $m=p^{k-1}$, and
$M=R^n$, $F=Q^n$ with positive integer bases. Then
\begin{align}
 c_{\mathrm{bad}}(n)&\le\frac23R^m<\frac23c^{2/p},
 \label{ppg-modn}\\
 27c_{\mathrm{bad}}(6n)^2&\le16R^{2m+1},
 \label{ppg-mod6n}\\
 \sum_{\substack{q\mid c\\q\equiv1\bmod6p^k}}v_q(c)\log q
 &>\left(1-\frac2p-\frac1{p^k}\right)\log c
       +\log\frac{3\sqrt3}{4}.
 \label{ppg-mass}
\end{align}
In particular, an actual prime divisor of $c$ is $1\bmod6p^k$,
and is at least $6p^k+1$.
\end{corollary}
\begin{proof}
Here $\Phi_n(1)=p$, $\kappa=1$ and $n-\varphi(n)=m$.
Substitution in the general support budgets proves all three estimates.
The exponent in \eqref{ppg-mass} is positive, and
$3\sqrt3/4>1$, so the good part exceeds one.

There is also a direct allocation proof that makes the $p$-exception
transparent. Set
\[
 D=Q^m-R^{2m},\qquad
 S=\Phi_n(Q,R^2)
  =\sum_{j=0}^{p-1}(Q^m)^{p-1-j}(R^{2m})^j.
\]
Then $DS=ab c^2$. At $q\ne p$ dividing $S$, the order of $Q/R^2$
is exactly $n$: its $n$-th power is one, its $m$-th power is not one,
and every proper divisor of $n$ divides $m$. At $p$, the binomial
argument gives $p\mid S$ exactly when $p\mid D$, and then $v_p(S)=1$.
Consequently $(ab c^2)_{\mathrm{bad}}(n)\mid pD$. Since
$S\ge pR^{2m(p-1)}$ and $ab c^2\le4R^{2n}/9$,
\[
 1\le D\le\frac{4R^{2m}}{9p},\qquad
 (ab c^2)_{\mathrm{bad}}(n)\le\frac49R^{2m}.
\]
This proves \eqref{ppg-modn} with full valuations. The actual input-arm
allocation for the remaining inert primes gives $3J^2\le4R$ and
$c_{\mathrm{bad}}(6n)=c_{\mathrm{bad}}(n)J$, proving
\eqref{ppg-mod6n} again.
\end{proof}

At $k=1$ this recovers the split progression bound at $6p$. For fixed
$p$ and growing $k$, the guaranteed fraction tends to $1-2/p$; it
does not tend to one merely because the modulus grows. Any common
prime-power divisor $p^k\mid\gcd(h,g)$ of unequal pure exponents
$M=A^h$, $F=B^g$ gives the same result by replacing their bases.

\begin{theorem}[The entire common exponent controls the extraction root]
\label{ppg-full-exponent-growth}
For any positive integer $n$, simultaneous representations $M=R^n$,
$F=Q^n$ with positive integer bases satisfy
\[
 9n(Q-R^2)\le4R^2,\qquad R^2\ge\frac{9n}{4},\qquad
 \log c>\frac n4\log\frac{9n}{4}.
\]
\end{theorem}
\begin{proof}
The inequality $F>M^2$ gives the positive integer $D_0=Q-R^2$.
The full geometric sum has $n$ terms, each at least $R^{2n-2}$,
and its product with $D_0$ is $ab c^2\le4R^{2n}/9$.
Cancellation gives the first inequality. Use $D_0\ge1$ for the
second, and $R^n=M<c^2$ for the logarithmic conclusion.
This includes $n=1$ and needs no prime-order assertion.
\end{proof}

This is a lower bound on the height of the actual seed. It is not an
upper bound, and does not contradict the existence of a moving family.

\paragraph{Exact scope of the fifty-four new Lean declarations.}
The signed-arm module proves thirty-two statements. It defines the actual
normalized Eisenstein power, proves boundary-arm divisibility, reconstructs
the three integer quotients, proves their additive relation and norm identity,
and derives pairwise coprimality with output primitivity explicitly supplied.
It also proves the signed linear budget, its two-arm specialization, finite
nonnegative integer-weight depth bounds, and the exact overlap-cost inequality.
The finite weights are not silently identified with real logarithms.

Two common-exponent modules prove fifteen statements. Besides the actual
polynomial certificates, they prove the geometric-sum identity and lower
bound for all relevant integer bases. In particular, the integer conclusions
of Theorem~\ref{ppg-full-exponent-growth} follow directly from actual norm
representations; the sum lower bound is proved rather than postulated.
The separate bad-part interface retains its explicit antecedent.

The support-assembly module proves seven statements. For every odd integer
$n$, the rank premises $n\mid q-1$, $d\mid n$, $d\mid q+1$ and
$d>0$ give $d=1$: they force $d\mid2$, and $d=2$ contradicts oddness.
The module proves the two actual input-arm height inequalities using
\[
 4N(x,y)-3x^2=(x+2y)^2,\qquad
 4N(x,y)-3(x+y)^2=(x-y)^2.
\]
It combines explicit budgets $9B^2\le4U$ and $3J^2\le4R$ into
$27(BJ)^2\le16UR$, including its power form. Finally it verifies the
integer core of the totient support-existence step. For $R\ge1$,
$0<K\le N$, $9N\le4R^2$ and integer $\rho\ge2$, the inequality
$16KG^2>27R^\rho$ forces $G^2>1$.
The correspondence of these parameters with actual prime parts remains
the independently reviewed ordinary input.

All fifty-four new declarations and fifty-four unchanged dependency
declarations pass a fresh scoped Lake build under Lean~4.32.0 with warnings
treated as errors. Every theorem's axiom list is checked; only
\texttt{propext}, \texttt{Classical.choice} and \texttt{Quot.sound} occur.
Three exact finite supplements are independently replayed: $113$ actual
signed-arm rows and $678$ rational inequalities; $116$ cyclotomic instances
and $3931$ primitive seed identities; and $14$ fully factored general-index
cyclotomic values with four additional exceptional-prime valuations.
These are finite checks, not an infinite family or a substitute for the
ordinary universal proofs.

The full prime-order and valuation allocation, unique factorization, real
logarithmic estimates, and arbitrary-root signed-tail membership are outside
this continuation's Lean scope. The complementary high multiplicities and
the absence of a suitable common exponent remain live bottlenecks. Neither
standard ABC nor its negation is asserted by any of these results.
